\documentclass[11pt]{article}
\usepackage[T1]{fontenc}
\usepackage{lmodern,amsmath,amssymb,amsthm,booktabs,xcolor,microtype}
\usepackage[letterpaper,margin=1in,top=.87in,bottom=.9in]{geometry}
\usepackage{fancyhdr}
\definecolor{ink}{HTML}{17384A}
\definecolor{link}{HTML}{145E68}
\usepackage[colorlinks=true,urlcolor=link,citecolor=link,linkcolor=link]{hyperref}
\hypersetup{pdftitle={Sparse equations for singular moduli and optimal tests for symmetric-power zeros}}
\newtheorem{theorem}{Theorem}[section]
\newtheorem{proposition}[theorem]{Proposition}
\newtheorem{lemma}[theorem]{Lemma}
\newtheorem{corollary}[theorem]{Corollary}
\pagestyle{fancy}\fancyhf{}\renewcommand{\headrulewidth}{.25pt}
\fancyhead[L]{\small\sffamily\color{ink}Singular moduli and symmetric-power zeros}
\fancyhead[R]{\small\sffamily\color{ink}Research note}
\fancyfoot[C]{\small\thepage}\setlength{\headheight}{15pt}
\setlength{\emergencystretch}{2.5em}\raggedbottom
\newcommand{\repo}{https://github.com/lsernau/grh-lean/blob/72fead102f663d5c7a407964a13dc0cf9d977207}
\begin{document}
\thispagestyle{plain}
\noindent{\small\sffamily\color{link}NUMBER THEORY\hfill 25 NOVEMBER 2026}\par\vspace{10pt}
{\LARGE\bfseries\color{ink}Sparse equations for singular moduli and optimal tests for symmetric-power zeros\par}\vspace{9pt}
\noindent\textbf{Abstract.}
For more than $99.9818260107372245\%$ of prime conductors in the Gaussian field, in lower relative density, a degree-$h$ singular modulus requires exactly $h+1$ terms in every sparsest rational equation. The exponents may be arbitrarily large. We improve the earlier $61.33\%$ lower bound by putting all conjugates into common multiplicative disks at $2$, $3$ and $7$ using isogenies of elliptic curves. The result is independent of RH. A separate transfer of the one-sided argument in \emph{grh-lean} gives an extremal problem for nonnegative weights on the normalized Hecke traces of Ramanujan's $\Delta$. We solve the optimization with the two leading coefficients prescribed. A single upper bound for the resulting weighted prime sum is equivalent to simultaneous RH for the symmetric powers it tests. Proving that arithmetic upper bound remains open.

\section{Which parts of the Lean proof transfer?}
The public statements in \texttt{grh/src/Statements.lean} and the five theorems in \texttt{Main.lean} concern the Riemann zeta function and Dirichlet $L$-functions over $\mathbb Q$. This note uses the readable edition at revision \texttt{72fead102f66}. The repository records kernel checks and its axiom set. Our investigation inspected the source and its declarations. We did not rebuild the Lean projects or formalize the results below. The new proofs use classical mathematics and exact rational computations.\cite{source}

\paragraph{The arithmetic input.}
The argument obtains an upper bound for a smoothed sum over prime powers $n\equiv1\pmod q$. Its coefficients are $\Lambda(n)(1-\chi_K(n))$ for an auxiliary quadratic character. On the geometric side, short integral forms bound a covolume from above, a spanning theorem controls how much of the space they occupy, and determinants of chains of theta lifts supply a lower bound. Comparing these for the same lattice bounds the excess divisibility. The formal junctions are \texttt{star\_inequality} and \texttt{Data.primeRemainder\_le}. The smoothing factor is $1/n-1/x$ for $n\le x$.\cite{source}

\paragraph{The analytic input.}
In \texttt{New/Reduction/Zeros/Criterion.lean}, the Mellin transform sees the orders of zeros of Dirichlet factors and their quadratic twists, with opposite signs. Given a proposed zero, an auxiliary twist is chosen to avoid zeros in a bounded-height region. Real residues then have the sign needed for Landau's theorem. Section~\ref{sec:analytic} replaces twist selection by weights with strictly positive coefficients on every symmetric-power character being tested. Section~\ref{sec:weights} computes the least main-term coefficient compatible with specified top coefficients and pointwise nonnegativity.

\paragraph{The earlier arithmetic question.}
Bilu, Luca and Pizarro-Madariaga showed that Dirichlet GRH excludes trinomials for singular moduli of degree at least three.\cite{blp} The incoming ProveIt report goes further and asks how often the least possible number of terms is the maximum $h+1$ for orders of prime conductor in $\mathbb Q(i)$. Section~\ref{sec:gaussian} advances this question through a stronger local root count. The connection to \emph{grh-lean} here is the research problem reached through the previous reports. The symmetric-power calculation directly adapts the repository's analytic argument.\cite{cmprior,deficit}

\section{Full support for Gaussian prime conductors}
\label{sec:gaussian}

For a singular modulus $j$ of degree $h$, define $s(j)$ as the minimum
number of nonzero coefficients in a nonzero polynomial $P\in\mathbb Q[X]$
with $P(j)=0$. The minimal polynomial gives $s(j)\le h+1$. Equality
means that \emph{every} $h$ distinct nonnegative powers of $j$ form a
basis of $\mathbb Q(j)$. There is no bound on the exponents here.

\medskip\noindent\textbf{Theorem.}
For a rational prime $\ell$, let $j_\ell$ have endomorphism order
$\mathbb Z+\ell\mathbb Z[i]$, of discriminant $-4\ell^2$. The set of
primes with $s(j_\ell)=[\mathbb Q(j_\ell):\mathbb Q]+1$ has lower
relative natural density greater than $0.999818260107372245$.
The proof is independent of RH and GRH. The exact density of the
full-support set remains open.

\paragraph{The earlier problem.}
The ProveIt report \emph{Five-term equations for singular moduli},
in its sections on prime conductors and further questions, derives a
lower density greater than $0.6133$ and asks for the true density.
Its local calculation puts all conjugates in one residue class modulo
$7$. We use good integral models of the curves at $3$ and $7$ and an
explicit isogeny calculation at $2$ to obtain three multiplicative disks. Then we use all factorial comparisons across a single
cut. The finite density computations below concern the resulting
sufficient conditions. The
\href{https://github.com/VladimirReshetnikov/ProveIt/blob/6fe2c1bd450aa7f19eab17fb9a6e76170b4fce00/docs/incoming/grh_sparse_singular_moduli.zip}{prior archive}
contains \texttt{local.tex}, \texttt{prime-conductors.tex} and
\texttt{closing.tex}. For the original trinomial problem see
\href{https://arxiv.org/abs/2003.06547}{Bilu--Luca--Pizarro-Madariaga}.

\paragraph{The common disks at odd primes.}
We may discard $\ell=2,3,7$. Write $\sigma=\chi_{-4}(\ell)=\pm1$.
The order class number formula gives $h=(\ell-\sigma)/2$, since
$\mathbb Z[i]$ has class number one and the unit index is two.
Every conjugate is the invariant of a cyclic degree-$\ell$ quotient
of $E_0:y^2=x^3-x$. Indeed its CM lattice corresponds to an invertible
ideal $I$ of the order. Passing from $I$ to $I\mathbb Z[i]$ gives index
$\ell$, as follows locally from invertibility. The extended lattice
is principal, so the resulting curve has invariant $1728$. Taking
the dual isogeny gives the desired quotient. For background on this
lattice description see
\href{https://arxiv.org/abs/1208.5370}{Sutherland, \emph{Isogeny volcanoes}}.

At $p=3,7$, the curve $E_0$ has good reduction. In each case a finite
unramified extension contains all its $\ell$-torsion, since the torsion
of order prime to $p$ is finite \'etale on the good model. Over its
integer ring $\mathcal O$, quotienting by each cyclic subgroup gives
an elliptic curve with good reduction. Its reduction is supersingular
because it is isogenous to the supersingular reduction of $E_0$.
The latter assertion also follows from the Hasse invariant computation
below applied to $E_0$. The good-reduction facts can equivalently be
deduced from the N\'eron--Ogg--Shafarevich criterion.

At $p=7$ use a good integral short equation $y^2=x^3+Ax+B$.
The Hasse invariant is the coefficient of $x^6$ in $(x^3+Ax+B)^3$,
which is $3B$. Thus $7\mid B$. The denominator $4A^3+27B^2$ is a
unit, and
\[
 \frac{j}{1728}=\frac{4A^3}{4A^3+27B^2}
                 =1-\frac{27B^2}{4A^3+27B^2}\in 1+7^2\mathcal O .
\]
At $p=3$, complete the square in a good integral equation to get
$y^2=x^3+a_2x^2+a_4x+a_6$. The Hasse invariant is the coefficient of
$x^2$ in this cubic, so $3\mid a_2$. Translating the $x$ coordinate
by $-a_2/3$ gives an integral short equation with unit discriminant.
The same formula now gives $j/1728\in 1+3^3\mathcal O$. Division by
$1728$ is taken in the local field at $3$.

\paragraph{A common multiplicative disk at 2.}
Let $v(2)=1$. For any two cyclic order-$\ell$ subgroups $H,J$ of
$E_0$, with $\ell$ odd, we claim
\[
 v\big(j(E_0/H)/j(E_0/J)-1\big)\ge4.                 \tag{G0}
\]
All valuations may be taken in one finite extension of $\mathbb Q_2$.
The integral invariant $1728$ gives $E_0$ potentially good reduction.
Take a field of good reduction containing $E_0[3\ell]$. Write
$x=r+u^2X$ for the change from the original short equation to a good
integral Weierstrass equation with coordinate $X$. Its discriminant
is a unit, while the original discriminant is $64$, so $v(u^2)=1$.
Every nonzero odd-torsion point under consideration has integral $X$:
reduction is injective on prime-to-$2$ torsion, so its reduction is
an affine point of the good model. In particular for a fixed
nonzero $3$-torsion point $T$ and every nonzero $Q\in H\cup J$ we
have $v(x_Q-x_T)\ge1$.

The $3$-division polynomial gives $3x_T^4-6x_T^2-1=0$.
Its roots are units. For $y=5x_T^2-3$ this yields
$3y^2-12y-88=0$. The Newton polygon has its only slope $-3/2$,
so $v(y)=3/2$. Since $x_Q=x_T+2z$ with $z$ integral,
$v(x_Q^2-x_T^2)\ge2$. For $g(x)=x(5x^2-3)$ it follows that
\[
 v(g(x_T))=3/2,\qquad v(g(x_Q)-g(x_T))\ge2.
\]
V\'elu's quotient model is $Y^2=X^3+A_H X+B_H$, where
\[
 A_H=-1-10\!\sum_{Q\in(H\setminus\{0\})/\{\pm1\}}(3x_Q^2-1),
 \qquad
 B_H=-14\!\sum_{Q\in(H\setminus\{0\})/\{\pm1\}}g(x_Q).
\]
These are the short-equation specialization of
\href{https://ddd.uab.cat/pub/pubmat/pubmat_a2007vEXTRA/pubmat_a2007vextrap147.pdf}{Miret--Moreno--Rio,
\emph{Generalization of V\'elu's formulae for isogenies between elliptic curves},
pp.~149--151, especially Theorem~1}.
In particular $A_H$ is a unit congruent to $-1$ modulo $2$ and
$v(B_H)\ge5/2$. The $H$ and $J$ sums have the same number
$(\ell-1)/2$ of terms, all congruent to $g(x_T)$ modulo $4$ in
the second sum. Thus $v(B_H-B_J)\ge3$ and
$v(B_H+B_J)=v(B_H-B_J+2B_J)\ge3$.
This gives $v(B_H^2-B_J^2)\ge6$.

Put $s_H=27B_H^2/(4A_H^3)$. Then $v(s_H)\ge3$ and
$v(s_H-s_J)\ge4$: replacing $B_H^2$ by $B_J^2$ costs at least four
in valuation after division by $4$, while
$v(A_H^{-3}-A_J^{-3})\ge1$. The identity
$j(E_0/H)=1728/(1+s_H)$ proves (G0).
Fixing any one of the quotient invariants $j_*$ therefore puts all
$h$ conjugates in $j_*(1+16\mathcal O)$, where $\mathcal O$ is the
valuation ring. The center is allowed to be algebraic. The recurrence
and the root count below apply over this complete finite extension
with the same normalized valuation.

The depth four is sharp when the disk assertion ranges over all odd
prime conductors. For $\ell=3$ let $Y=j/1728-1$. Substitution of
$A=9-30x_T^2$, $B=-14x_T(5x_T^2-3)$ into the formula for $j$
and elimination by $3x_T^4-6x_T^2-1=0$ give the quadratic
$27Y^2-2399040Y-18593344$. Its discriminant has $2$-adic valuation
eight. The two roots differ by an element of valuation four, and
both $1+Y$ are units. The corresponding pair thus attains equality
in (G0). The archive reproduces the elimination and also checks
the corresponding degree-five isogenies in characteristic zero.


\paragraph{The factorial comparison.}
Use $(p,b)=(2,4),(3,3)$ or $(7,2)$ and $w(n)=bn-v_p(n!)$, with $w(0)=0$.
The sufficient condition at $p$ is
\[
       \max_{0\le u<h}w(u)<\min_{v\ge h}w(v).                 \tag{G1}
\]
To prove this, suppose that a polynomial with $k\le h$ terms and
distinct exponents $e_j$ vanishes at a conjugate. It vanishes at all
$h$ conjugates. Let $u_0$ be the common multiplicative center given by these disks,
namely $1728$ at $3,7$ and a chosen conjugate at $2$. Compose with
$X=u_0\exp(p^bY)$ on the closed unit disk of a finite local field
containing the center and all conjugates. The logarithm gives a distinct point
$Y$ for each root. The coefficient at degree $n$ is $p^{bn}A_n/n!$,
where $A_n=\sum_j c_ju_0^{e_j}e_j^n$. A Vandermonde determinant
gives $M=\max_{n<k}|A_n|_p>0$. The monic integral recurrence with
characteristic polynomial $\prod_j(T-e_j)$ implies $|A_n|_p\le M$
for all $n$. Under (G1), all coefficients from $n=h$ onward have
smaller norm than the coefficient at an index below $k$ attaining
$M$. The series converges, since $b>1/(p-1)$.
Strassmann's theorem bounds the number of its roots by $h-1$, a
contradiction. This is the disk argument of the prior report applied
to the deeper common disks. For context on fewnomial root counts see
\href{https://arxiv.org/abs/1008.4808}{Avenda\~no--Krick}.

\paragraph{Counting the prime conductors.}
Taking differences $w(v)-w(u)$ shows that (G1) is equivalent to
\[
        \sum_{t\in I}v_p(t)<b|I|
        \quad\text{for every integer interval }I\subset\mathbb Z_{>0}
        \text{ containing }h.                                \tag{G2}
\]
For lengths $r\le R$, omitting the finitely many $h<R$, the test
depends only on $h\bmod p^{bR}$. Fix $\sigma$. For an odd $p$, the admissible classes of $h$ modulo $p^a$
are those with $p\nmid 2h+\sigma$. For $p=2$, they are exactly the
even classes: $\ell=2h+\sigma\equiv\sigma\pmod4$,
and all such classes correspond to invertible classes of $\ell$
modulo $2^{a+1}$ with that sign. In all three cases each admissible class
has conditional density $1/\varphi(p^a)$, by the prime number theorem
in arithmetic progressions. Reflection $h\mapsto-h$ and reversal of
intervals show that the finite density $d_{p,R}$ is the same for
either $\sigma$.

Here is a uniform bound passing from finite lengths to the infinite
test. If a length-$r$ interval violates (G2) and its greatest
valuation is $m$, counting multiples of $p^a$ for $1\le a\le m$
bounds the sum of valuations by $r/(p-1)+m$. Consequently some
$h+j$, with $|j|\le r-1$, is divisible by
\[
      q_r=p^{\lceil (b-1/(p-1))r\rceil}.
\]
For fixed $r$ the conditional upper density of this event is at most
$u_r=(2r-1)/\varphi(q_r)$. To justify summing the bound, write
$X$ for the prime cutoff. For $q_r\le\sqrt X$ the Brun--Titchmarsh
inequality bounds the possible primes in each relevant invertible residue
class by $O(X/(\varphi(q_r)\log X))$, with an absolute constant.
The modulus is $q_r$ for odd $p$ after dropping the condition modulo
$4$, and $2q_r$ for $p=2$, because $\ell=2h+\sigma$. In the latter
case $2q_r\le 2\sqrt X$ still gives $\log(X/(2q_r))\gg\log X$. Since only $\ell=p$
can occupy a noninvertible class, dropping that prime and dividing
by the count of primes of each sign gives a uniform tail bounded by
an absolute multiple of $\sum_{r>T}u_r$. For $q_r>\sqrt X$, a positive
multiple of $p^{\lfloor\log_p\sqrt X\rfloor}$ occurs among the $h+j$.
The inequality $(b-1/(p-1))r\le\log_p(h+r)$ forces $r=O(\log X)$:
when $r>h$ the inequality fails since $h+r<2r$ and
$2r\le 2^r<p^{(b-1/(p-1))r}$. Otherwise $h+r\ll X$.
There are $O(\sqrt X\log X)$ possibilities for a multiple and an
offset, negligible compared with the number of primes. Taking limits
at fixed $T$ and then $T\to\infty$ proves existence of the test
density $d_p$ and
\[
         0\le d_{p,R}-d_p\le\sum_{r>R}u_r.                  \tag{G3}
\]
The sieve inequality used here is due to
\href{https://doi.org/10.1112/S0025579300004708}{Montgomery--Vaughan}.

\paragraph{Exact certificate.}
The verifier starts with the $p-1$ admissible classes modulo $p$.
In a class modulo $p^a$, it evaluates all relevant valuations truncated
at $a$. If any interval sum already reaches $br$ it discards the
class. If no valuation is truncated it accepts the class. Otherwise
it visits all $p$ lifts. The process terminates by depth $bR$, since
a truncated term at that depth already violates a length at most $R$
containing that term and the center. To bound (G3), group its terms
in successive blocks of length $p-1$. For $r\ge R+1$,
\[
 \frac{u_{r+p-1}}{u_r}=
 \frac{2r+2p-3}{2r-1}\,p^{-(b(p-1)-1)}
\]
is decreasing. A geometric sum bounds the whole tail using its first
block. All operations in \texttt{check\_density.py} use rational
arithmetic for the bounds.
\[
\begin{array}{c|c|c|c}
p & R & d_{p,R} & \text{tail upper bound}\\ \hline
2&25&65963232382760376856543/75557863725914323419136 &3.879\cdot10^{-22}\\
3&16&68736524642595367189/72945992754341572806 &1.694\cdot10^{-19}\\
7&10&9528160153068835/9770481587462694 &5.089\cdot10^{-17}
\end{array}
\]
At fixed sign, the three finite congruence tests are independent by the
Chinese remainder theorem and the prime number theorem in progressions.
Passing to the limits using (G3), the density of passing at least one test
is $d=1-(1-d_2)(1-d_3)(1-d_7)$. It is the same for both signs and satisfies
\[
 0.9998182601073722452943 < d < 0.9998182601073722456678 .
\]
Every prime counted by $d$ gives full support. Determining what occurs
for the remaining conductors would require information beyond these
three disk tests.

\paragraph{Individual conductors.}
For $\ell=197$ the degree is $98$. The one-term interval $\{98\}$
fails (G2) at $7$, while it passes at $3$. At $3$ every interval of
length $r\ge2$ through $98$ has valuation sum at most
$r/2+\log_3(98+r)<3r$. The last inequality holds at $r=2$ and
persists since $3^{5r/2}$ grows faster than $98+r$ at every step.
Thus any $98$ distinct nonnegative powers of $j_{197}$ form a basis.
For $\ell=107$ the degree is $54$. Here $3^3\mid54$ rules out (G2)
at $3$, while $7\nmid54$. For $r\ge2$ one has
$r/6+\log_7(54+r)<2r$, again by checking $r=2$ and increasing $r$.
Every $54$ distinct powers of $j_{107}$ therefore form a basis.

Conductor $15877$ exhibits the contribution at $2$. It is prime
(by trial division through $126$) and $h=7938=2\cdot 3^4\cdot7^2$.
Both odd-prime tests fail on the singleton $\{h\}$. The largest
sums of $2$-adic valuations of intervals through $7938$ of lengths
$1,2,3,4$ are respectively $1,1,9,9$, each smaller than four times
the corresponding length. For every $r\ge5$, the sum is at most
$r+\log_2(7938+r)<4r$, since $7938+5<2^{15}$ and the inequality
persists on increasing $r$. Thus any $7938$ distinct nonnegative
powers of this singular modulus form a basis.


\paragraph{Finite verification and a remaining example.}
Among the $9{,}589$ primes $\ell\le 100{,}000$ after excluding $2,3,7$, exactly $9{,}587$ pass at least one test. The two odd primes alone certify $9{,}575$. The first conductor missed by all three is $\ell=42{,}337$, with $h=21{,}168=16\cdot27\cdot49$. The singleton $\{h\}$ defeats all three tests. The other is $84{,}673$ with $h=42{,}336$. The exact least support in these two cases is undetermined here. These counts come from a separate direct implementation of (G1). For any given $h$, it can stop as soon as $n(b-1/(p-1))>bh$, since this lower bound for $w(n)$ then exceeds every $w(u)$ for $u<h$. The infinite density estimate comes from (G3).

\section{The least possible mean for a positive symmetric-power test}\label{sec:weights}
Write $t=2\cos\theta$ and
\[
 C_j(t)=U_j(t/2)=\frac{\sin((j+1)\theta)}{\sin\theta},
 \qquad d\mu(t)=\frac{\sqrt{4-t^2}}{2\pi}\,dt,\quad -2\le t\le2.
\]
At the endpoints use the continuous values. The $C_j$ are monic and
orthonormal for the Sato--Tate probability measure $\mu$. We use
$C_{-1}=0$, $C_0=1$, and $tC_j=C_{j+1}+C_{j-1}$.
Thus the coefficient of $C_0$ is the Haar mean $\int P\,d\mu$,
and the coefficient of $C_j$ is the multiplicity weight assigned to
the $j$th symmetric power. Since $C_j$ has parity $j$, the two
highest character coefficients of a polynomial are also its two
highest ordinary coefficients. The character product identity is
\[
 C_i C_j = \sum_{h=0}^{\min(i,j)} C_{i+j-2h}.
\]

\subsection{A quadrature certificate}
Fix $k\ge1$, let $u=(k+1)/k$, and define $Q_a=C_k+aC_{k-1}$
for $0\le a\le u$. Write its zeros as $z_1<\cdots<z_k$.
They are all interior to $[-2,2]$ except that $z_1=-2$ when
$a=u$. To see this for $k\ge2$, evaluate $Q_a$ at the zeros of
$C_{k-1}$, where $Q_a=C_k$ has alternating signs by the sine formula.
The values $Q_a(2)=k+1+ak$ and $(-1)^k Q_a(-2)=k+1-ak$ supply
the two outer roots, counting $-2$ when equality holds. The degree
gives simplicity. For $k=1$ the zero is $-a$.

For the cardinal interpolation polynomials $\ell_i$ at the $z_i$,
put $\lambda_i=\int \ell_i\,d\mu$. For $\deg V\le2k-2$ division
by $Q_a$ gives a quotient of degree at most $k-2$, orthogonal to
$Q_a$, and a remainder equal to $\sum_i V(z_i)\ell_i$. Therefore
\[
 \int V\,d\mu=\sum_i\lambda_i V(z_i),
 \qquad \lambda_i=\int\ell_i^2\,d\mu>0.
\]
The second identity follows by taking $V=\ell_i^2$.
This elementary form of quasi-Gaussian quadrature is what makes the
optimization certify its own lower bound.

\subsection{Odd degree}
\begin{proposition}\label{prop:oddfront}
Among monic polynomials of degree $2k-1$ nonnegative on $[-2,2]$,
the least mean is $(k+1)/k$. The unique minimizer is
\[
 R_k(t)=\frac{Q_u(t)^2}{t+2}.
\]
Every character coefficient of $R_k$ is strictly positive.
\end{proposition}
\begin{proof}
Since $Q_u(-2)=0$, $R_k$ is a monic polynomial, nonnegative everywhere
on the interval, with a simple endpoint zero and $k-1$ double
interior zeros. If $P$ is a competitor, $P-R_k$ has degree at most
$2k-2$, so the rule at $a=u$ gives
$\int(P-R_k)\,d\mu=\sum_i\lambda_i P(z_i)\ge0$.
Equality forces the endpoint zero and every interior double zero.
Together with monicity they determine $R_k$. Writing $[C_j]R_k$ for
its character coefficients, their values are
\begin{align*}
 [C_{2j}]R_k&=\frac{k(k+1)-j(j+1)}{k^2}
                  &&(0\le j\le k-1),\\
 [C_{2j-1}]R_k&=\frac{j^2}{k^2}
                  &&(1\le j\le k).
\end{align*}
For a direct check multiply these coefficients by $t+2$: the coefficient
at $j$ becomes $2r_j+r_{j-1}+r_{j+1}$, with missing terms set to zero.
This is $1+u^2$ at even $j<2k$, $2u$ at odd $j<2k$, and $1$
at $2k$, the coefficients of $Q_u^2$ by the product identity.
In particular $r_0=u$ and every displayed coefficient is positive.
\end{proof}

\subsection{Even degree with specified weight on the preceding power}
For even degree, minimizing the mean with only the top coefficient
specified yields $C_k^2$, which omits every odd symmetric power.
Prescribing the coefficient $b$ of $C_{2k-1}$ quantifies the cost of
including the preceding power. The other coefficients are free in the
minimization.

\begin{theorem}\label{thm:evenfront}
Let $P$ be monic of degree $2k$, nonnegative on $[-2,2]$, with
$[C_{2k-1}]P=b\ge0$. The least possible mean is
\[
 M_k(b)=
 \begin{cases}
  1+b^2/4,&0\le b\le 2u,\\
  1-u^2+ub,&b\ge 2u .
 \end{cases}
\]
The unique minimizing polynomial is
\[
 P_{k,b}=
 \begin{cases}
  Q_{b/2}^2,&0\le b\le2u,\\
  Q_u^2+(b-2u)R_k,&b\ge2u .
 \end{cases}
\]
For $b>0$ every character coefficient of $P_{k,b}$ is strictly
positive. Thus imposing this additional requirement preserves the
minimum and its attainment.
\end{theorem}
\begin{proof}
For $0\le a\le u$, subtract $Q_a^2+(b-2a)C_{k-1}Q_a$ from
$P$ to cancel the two highest terms. Applying the quadrature rule to
the difference and using orthonormality gives the exact identity
\[
 \int P\,d\mu=1-a^2+ab+\sum_{i=1}^{k}\lambda_i P(z_i).
\]
Maximizing the quadratic lower bound over $[0,u]$ selects
$a=\min(b/2,u)$ and gives $M_k(b)$. In the first range $Q_a^2$
attains equality. For $b\ge2u$ the displayed sum of $Q_u^2$ and
$(b-2u)R_k$ is nonnegative, has the prescribed top two coefficients,
and vanishes at all quadrature nodes, so also attains equality.
In the interior-node case equality forces $k$ double zeros and hence
the monic square. At $a=u$ equality forces divisibility by $R_k$.
The remaining monic linear factor is uniquely specified by the
coefficient $b$. This includes the common boundary $b=2u$.

The product identity gives $[C_{2k}]Q_a^2=1$, all lower even
coefficients $1+a^2$, and all odd coefficients $2a$.
Proposition~\ref{prop:oddfront} supplies the additional coefficients
in the second range.
\end{proof}

\noindent
For $0<b\le2u$ the same $M_k(b)$ is the sharp minimum when
\emph{every} odd coefficient is required to be at least $b$.
Indeed a competitor has top odd coefficient $b'\ge b$ and
$M_k(b')\ge M_k(b)$, while $Q_{b/2}^2$ has all odd coefficients
exactly $b$. For $b>2u$ Theorem~\ref{thm:evenfront} prescribes the
top odd coefficient, and its remaining odd coefficients are given by
the formulas above. Without a prescribed positive lower bound on any
odd coefficient, the infimum among strictly positive character
coefficients is $1$, approached as $b\downarrow0$. The unique
unconstrained minimizer $C_k^2$ shows this infimum is unattained
under strict positivity.

\paragraph{Relation to the sine-polynomial literature.}
Multiplication by $\sin\theta$ sends $P(2\cos\theta)=\sum c_j C_j(2\cos\theta)$
to the sine polynomial $\sum c_j\sin((j+1)\theta)$ nonnegative on
$[0,\pi]$. Hence a result about the top coefficient after normalization
by the first sine coefficient is directly relevant here. Rogosinski
and Szeg\H{o} introduced this extremal setting~\cite{rs}.
Dimitrov and Merlo~\cite{dm}, Theorems~1 and~2, give the individual
top-coefficient extremizers, including the odd-degree minimum above
and the even-degree example $C_k^2$. Their next-to-top coefficient
extremizer for odd sine degree is the mean-normalized case $b=2$ of
Theorem~\ref{thm:evenfront}. The joint formula fixes both highest
coefficients and describes the transition at $2+2/k$. Royster and Suffridge also determined the low-degree coefficient region for
typically real cubics, which under the same sine substitution covers
$k=1$ here.\cite{roysuff} We give the joint formula in all degrees and
make no assertion of bibliographic priority.

\section{A one-sided criterion for the symmetric powers of \texorpdfstring{$\Delta$}{Delta}}\label{sec:analytic}

Write $\Delta(z)=\sum_{n\ge1}\tau(n)e^{2\pi i nz}$ for Ramanujan's weight-twelve level-one newform. At each prime $p$, let
\[
 t_{p,m}=\alpha_p^m+\alpha_p^{-m},\qquad
 \alpha_p+\alpha_p^{-1}=\tau(p)p^{-11/2},\qquad |\alpha_p|=1 .
\]
The last equality is Deligne's bound. The level is one, so every finite Euler factor is unramified. Put $L_j(s)=L(s,\operatorname{Sym}^j\Delta)$ in unitary normalization, with $L_0=\zeta$. For $j\ge1$ symmetric-power automorphy and cuspidality are available for this non-CM form by Newton and Thorne \cite{nt}. We use the standard continuation, zero-free line $\Re s=1$, functional equation about $1/2$, and fixed-degree explicit formula for these functions \cite{nt,rtc}.

Let $P=\sum_{j=0}^{d}a_j C_j$ with $a_j>0$ for $0\le j\le d$, write $A=a_0$, and set
\begin{align*}
 S_P(x)&=\sum_{p^m\le x}(\log p)P(t_{p,m})
             \left(\frac1{p^m}-\frac1x\right),\\
 D_P(s)&=-\sum_{j=0}^{d}a_j\frac{L'_j}{L_j}(s),\qquad
 c_P=\lim_{s\to1}\left(D_P(s)-\frac{A}{s-1}\right)-A,\\
 R_P(x)&=S_P(x)-A\log x-c_P.
\end{align*}
Here $m$ ranges over positive integers. The identity
$C_j(t_{p,m})=\sum_{h=0}^j\alpha_p^{m(j-2h)}$
makes $D_P$ the Dirichlet series with coefficient $(\log p)P(t_{p,m})$ at $p^m$. One copy of the principal pole appears with weight $A$. The Haar extremizer supplies a $P$ which is also pointwise nonnegative, an additional property relevant to its arithmetic realization.

\begin{theorem}
For each fixed $P$ as above, the following are equivalent.
\begin{enumerate}
\item Every nontrivial zero of every $L_j$, $0\le j\le d$, lies on $\Re s=1/2$.
\item For every $0<\varepsilon<1/2$ there exist $K>0$ and $X\ge2$ such that
\[
 R_P(x)\le Kx^{-1/2+\varepsilon}\qquad(x\ge X).
\]
\end{enumerate}
The constants may depend on $P$ and $\varepsilon$.
\end{theorem}
\begin{proof}
For $\Re s>1$, termwise integration gives
\[
 \int_1^\infty R_P(x)x^{-s}\,dx
 =\frac{D_P(s)}{s(s-1)}-\frac{A}{(s-1)^2}-\frac{c_P}{s-1}.
\]
Our choice of $c_P$ removes both negative Laurent coefficients at $1$. At an off-line zero $\rho$ in the open strip, the residue of the right side is
\[
 \frac{\sum_j a_j\operatorname{ord}_{\rho}L_j}{\rho(1-\rho)}.
\]
It is nonzero at every such zero and strictly positive at a real zero. Common zeros contribute their orders additively.

Assume (2) and fix $\beta=1/2+\varepsilon$. For $x\ge X$, the function
$Kx^{\beta-1}-R_P(x)$ is nonnegative. Its Mellin integral with measure $x^{-s}\,dx$ converges for $\Re s>1$ by boundedness of $P$ and the elementary Chebyshev estimate. It has a meromorphic continuation in $\Re s>\beta$ by the displayed formula. If its real abscissa of convergence exceeded $\beta$, Landau's theorem for the transform of a nonnegative function would give a singularity there. The only possible singularities in this half-plane are simple poles from zeros of the $L_j$. At a real pole the residue of this transform is strictly negative, which would make the positive integral tend to $-\infty$ from the right. Removing the initial interval $[1,X]$ contributes an entire term. Thus the integral converges and is holomorphic on $\Re s>\beta$. Every zero there would produce a nonzero residue. Letting $\beta$ decrease to $1/2$ and applying the functional equations proves (1). This is the one-sided Mellin argument used in the repository \cite{source}; related transfers to weighted prime sums are discussed in the incoming reports \cite{deficit}.

Conversely, write $\psi_P(y)=\sum_{p^m\le y}\log p\,P(t_{p,m})$. Under (1), the fixed-degree explicit formula gives
$\psi_P(y)=Ay+O_P(\sqrt y\log^2(2y))$.
Indeed each zero term in the truncated formula is $O(\sqrt y/|\rho|)$, and the number of zeros up to height $T$ is $O_j(T\log(2T))$. Equivalently one may apply the fixed symmetric-power prime bound under GRH, then restore the higher prime powers by Chebyshev's estimate \cite{rtc}. Since $S_P(x)=\int_1^x\psi_P(y)y^{-2}\,dy$, the integral of $(\psi_P(y)-Ay)y^{-2}$ over $[1,\infty)$ converges and its tail is $O_P(x^{-1/2}\log^2(2x))$. The value of the integral is $c_P$, as follows by expanding $D_P(s)=s\int_1^\infty\psi_P(y)y^{-s-1}\,dy$ at $1$. This proves (2).
\end{proof}

\begin{proposition}
Let $R_P^{(1)}(x)$ use only $m=1$ in $S_P(x)$, and replace $c_P$ by
\[
 c_P^{(1)}=c_P-\sum_{p,\ m\ge2}\frac{\log p}{p^m}P(t_{p,m}).
\]
Then, unconditionally for this form,
\[
 R_P^{(1)}(x)-R_P(x)=\frac{2B_P+o(1)}{\sqrt{x}},
 \qquad B_P=\sum_{j=0}^{d}(-1)^j a_j.
\]
For the odd-degree extremizer $R_r$ of degree $2r-1$, this gives
$B_P=(r+1)(2r+1)/(6r)$. The one-sided criterion can therefore be stated using primes alone.
\end{proposition}
\begin{proof}
The correction to the centering constant is absolutely convergent. Put
$H(y)=\sum_{p^m\le y,\ m\ge2}\log p\,P(t_{p,m})$.
The terms $m\ge3$ are $O_P(y^{1/3}\log^2(2y))$ by the elementary Chebyshev estimate. For the squares, $t_{p,2}=t_{p,1}^2-2$. Sato--Tate for $\Delta$, the prime number theorem and partial summation give $H(y)=B_P\sqrt y+o(\sqrt y)$, with $B_P=\int P(t^2-2)\,d\mu(t)$. Indeed,
\[
 C_j(t^2-2)=\sum_{h=0}^j(-1)^h C_{2(j-h)}(t),
\]
by the character formula, so only the last term has nonzero Haar mean. Inserting the explicit coefficients of the extremizer and summing the two quadratic sequences yields the stated value. Partial summation gives
\[
 R_P^{(1)}(x)-R_P(x)=\frac{H(x)}{x}+\int_{(x,\infty)}\frac{dH(y)}{y}
   =\int_x^\infty\frac{H(y)}{y^2}\,dy,
\]
which proves the asymptotic, including its sign.
\end{proof}

For $d=1$, the pointwise weight is $P(t)=2+t$, and the criterion already tests the zeros of the degree-two $L$-function of $\Delta$ alongside those of $\zeta$. The missing arithmetic step is to prove the displayed upper bound for this trace-weighted remainder, at every sufficiently large $x$ for each $\varepsilon>0$. The repository's covolume comparison realizes a congruence and a quadratic-character weight. None of its short-form, spanning or determinant statements supplies the normalized trace $\tau(p)p^{-11/2}$. For higher $d$, the same question concerns the local weights $P(t_{p,m})$ constructed above.


\section{Further research questions}
\begin{enumerate}
\item Is the true density of full support equal to one? At the first missed conductor, $42{,}337$, all three first-difference tests fail at their central integer. A tied coefficient valuation leaves room for a nonzero generalized Vandermonde determinant. Its leading residue for the actual CM points is a more precise target than the valuations of the factorials alone.
\item Can the common-disk argument be extended to other CM isogeny classes? The proof singles out three requirements: a model with good reduction over an unramified field, a prime-to-residue-characteristic isogeny, and explicit control of the local coordinate of every quotient. The extra depth at $3$ comes from normalizing by $1728$ before taking logarithms.
\item Can a geometric divisibility inequality realize even the degree-one weight $2+\tau(p)/p^{11/2}$? The coefficient optimization controls the main term after fixing sensitivity to the highest symmetric powers. The needed global bound involves the actual Hecke traces and every sufficiently large scale. Improving an extremal polynomial alone does not supply that estimate.
\end{enumerate}

\begin{thebibliography}{9}\small\raggedright
\bibitem{source} L. Sernau and collaborators, \emph{GRH for Dirichlet $L$-functions, in Lean 4}. \href{https://github.com/lsernau/grh-lean/tree/72fead102f663d5c7a407964a13dc0cf9d977207}{Pinned source}, \href{\repo/grh/GUIDE.md}{reading guide}, \href{\repo/grh/src/Statements.lean}{public statements}, \href{\repo/grh/src/New/Reduction/MainInequality.lean}{covolume comparison}, and \href{\repo/grh/src/New/Reduction/Zeros/Criterion.lean}{zero criterion}.
\bibitem{cmprior} ProveIt, \emph{Five-term equations for singular moduli}, especially \texttt{local.tex}, \texttt{prime-conductors.tex}, and \texttt{closing.tex}. \href{https://github.com/VladimirReshetnikov/ProveIt/blob/6fe2c1bd450aa7f19eab17fb9a6e76170b4fce00/docs/incoming/grh_sparse_singular_moduli.zip}{Archived report}.
\bibitem{deficit} ProveIt, \emph{grh\_deficit\_transfer}, section on modular and symmetric-power RH. It gives the separate weight $r+1-C_r$ for a fixed $r$. \href{https://github.com/VladimirReshetnikov/ProveIt/blob/6fe2c1bd450aa7f19eab17fb9a6e76170b4fce00/docs/incoming/grh_deficit_transfer_2026-11-25.zip}{Archived report}.
\bibitem{blp} Y. Bilu, F. Luca and A. Pizarro-Madariaga, \emph{Trinomials, singular moduli and Riffaut's conjecture}, Theorem 1.2. \href{https://arxiv.org/abs/2003.06547}{arXiv:2003.06547}.
\bibitem{rs} W. W. Rogosinski and G. Szeg\H{o}, \emph{Extremum problems for non-negative sine polynomials}, Acta Sci. Math. (Szeged) 12, 112--124. \href{https://acta.bibl.u-szeged.hu/13628/1/math_012_pars_b_112-124.pdf}{Original paper}.
\bibitem{dm} D. K. Dimitrov and C. A. Merlo, \emph{Nonnegative trigonometric polynomials}, Theorems 1--2. \href{https://www.dcce.ibilce.unesp.br/~dimitrov/papers/ptpcapdf.pdf}{Authors' manuscript}.
\bibitem{roysuff} W. C. Royster and T. J. Suffridge, \emph{Typically real polynomials}, Publ. Math. Debrecen 17 (1970). \href{https://doi.org/10.5486/PMD.1970.17.1-4.35}{Publisher's version}.
\bibitem{nt} J. Newton and J. A. Thorne, \emph{Symmetric power functoriality for holomorphic modular forms}, I and II, Publ. Math. IH\'ES 134 (2021). \href{https://doi.org/10.1007/s10240-021-00127-3}{Part I} and \href{https://doi.org/10.1007/s10240-021-00126-4}{Part II}.
\bibitem{rtc} J. Rouse and J. Thorner, \emph{The explicit Sato--Tate conjecture and densities pertaining to Lehmer-type questions}, Trans. Amer. Math. Soc. 369 (2017), 3575--3604. \href{https://arxiv.org/abs/1305.5283}{Authors' version}.
\end{thebibliography}
\end{document}
